\chapter{ALIENS and UNDERGROUND BASES}\label{ch:underground}
\markboth{Aliens and Underground Bases}{Aliens and Underground Bases}

\chapquote{``I'm kinda sad, though. All these underground tunnels and lasers, but instead of having a rave we're using them for evil.''}{Gordon Freeman}

% ---------------------------------------------------------------
\section{The Greenbrier Resort}\label{sec:underground}

I want to begin this chapter with a real underground government facility, because the existence of
real ones is what makes the imaginary ones plausible, and because the Greenbrier is the perfect
case: it was secret, it was enormous, it was hidden under a hotel full of guests, and it stayed
secret for thirty years.

The Greenbrier is a luxury resort in White Sulphur Springs, West Virginia, about 250 miles from
Washington. Beginning in 1958 under the codename Project Greek Island, the Army Corps of Engineers
excavated a 112,\!000-square-foot bunker beneath the resort's West Virginia Wing, designed to house
the entire United States Congress --- both chambers, plus staff, roughly 1,\!100 people --- in the
event of nuclear war. It was completed in 1962, the year of the Cuban Missile Crisis, and remained
on 24-hour readiness until 1992.

\figwrap[l]{0.38}{The Greenbrier bunker blast door at White Sulphur Springs, concealed behind a wallpapered partition until 1992.}{https://commons.wikimedia.org/wiki/File:Greenbrier_bunker_blast_door.jpg}{fig:greenbrier}
The details are marvellous. Blast doors weighing 25 and 28 tonnes (see Figure~\ref{fig:greenbrier}), set in concrete walls up to five
feet thick, disguised behind ordinary wallpapered partitions. A dormitory with 18 bunkrooms. A
clinic with an operating theatre and a 12-bed intensive care unit. A crematorium. A decontamination
chamber. Water storage, diesel generators and three 25,\!000-gallon fuel tanks. A television studio
with a photographic backdrop of the Capitol, so broadcasts could appear to originate in Washington.
And the two largest halls --- the Governor's Hall and the Mountaineer Room --- were used
continuously as the resort's exhibition and conference space, with hotel guests holding trade shows
inside the House and Senate chambers for three decades without knowing it.

The cover was a government-contracted data facility, staffed by employees of a fictitious company,
``Forsythe Associates'', who maintained the bunker and serviced the hotel's televisions as a day
job. Supplies were rotated. Congressional leadership was briefed. And nobody talked, until Ted Gup
published an investigation in the \emph{Washington Post Magazine} on 31 May 1992, after which the
facility was decommissioned. You can now take the tour, which I recommend.

Why start here? Because the Greenbrier proves several things that matter for everything that
follows. Large secret facilities are genuinely built. They can be hidden in plain sight under
civilian cover. Secrecy can hold for thirty years across hundreds of participants. And ---
critically --- it ended with a journalist, a document trail, and a public tour. The secret was
kept, and then it was not, and the mechanism of its ending was ordinary reporting. Hold that
thought for Section~\ref{sec:dulce}.

% ---------------------------------------------------------------
\section{Cheyenne Mountain}\label{sec:cheyennemountain}

The other real one, and the model for every fictional underground base since.

\figwrap[r]{0.38}{A spring-mounted building inside the Cheyenne Mountain Complex, isolated from ground shock.}{https://commons.wikimedia.org/wiki/File:Cheyenne_Mountain_building_springs.jpg}{fig:cheyenne}
The Cheyenne Mountain Complex, near Colorado Springs, was excavated between 1961 and 1966 by
removing around 700,\!000 tonnes of granite, creating a grid of 15 three-storey buildings inside a
network of chambers roughly 600 metres beneath the summit. The buildings sit on more than a thousand steel coil springs (see Figure~\ref{fig:cheyenne}), each weighing about half a tonne, isolating them from ground shock ---
the complex was designed to survive a near miss by a multi-megaton weapon and continue operating.
Access is through a tunnel with two 25-tonne blast doors, and the tunnel runs through the mountain
rather than terminating, so that overpressure passes out the other side.

Inside: independent power generation, a six-million-gallon water reservoir in excavated rock,
NBC filtration, its own fire department and medical facility, and enough supplies for a sealed
crew. It was NORAD's command centre for missile warning and air sovereignty from 1966, and it is
the reason a generation grew up with a mental image of officers watching a big map.

In 2006 NORAD and USNORTHCOM operations moved to Peterson Space Force Base, and Cheyenne Mountain
was placed on warm standby as an alternate command centre. It remains active in that role, hosts
the Cheyenne Mountain Space Force Station, and its existence, dimensions, purpose and history are
completely public --- there are official photographs and published construction records.
You can see inside the complex here: \srclink{https://youtu.be/nVlVu397fNA}

Also real, also public: the Raven Rock Mountain Complex in Pennsylvania, ``Site R'', the Pentagon's
alternate; Mount Weather in Virginia, the FEMA special facility; the Kirtland Underground Munitions
and Maintenance Storage Complex; the Utah Data Center; and the Air Force Underground Test Area
facilities. The United States has built a great many holes in mountains and does not particularly
hide the fact.

Which brings me to the point of these two sections. When someone tells you the government has vast
underground facilities, the correct answer is yes, obviously, here are their names and
construction dates. The claim in Section~\ref{sec:dulce} is not that underground bases exist. It is a specific,
much stronger claim about what is in them.

% ---------------------------------------------------------------
\section{Dulce, NM}\label{sec:dulce}

And now the actual claim, which is the Dulce base, and I will lay it out fully and then explain why
it fails.

The story holds that beneath Archuleta Mesa near Dulce, New Mexico --- on Jicarilla Apache land ---
lies a joint human--alien facility of seven levels. Level 1 security and garages. Level 2 tube
shuttle station and offices. Level 3 government offices and laboratories. Level 4 mind-control and
aura research. Level 5 alien housing. Level 6, called ``Nightmare Hall'', containing genetic
experiments and multi-species hybrids. Level 7 holding thousands of humans in cold storage vats,
labelled and catalogued. The 1979 ``Dulce Wars'' are said to have been a firefight in which
sixty-six special forces personnel died.

\hi{The provenance is fully traceable, which is what makes this case instructive rather than merely
lurid.}

It originates with \textbf{Paul Bennewitz}, an Albuquerque electronics businessman who in 1979--80
became convinced he was intercepting alien communications from beneath Dulce, based on radio noise
and on lights over the mesa. He was located next to Kirtland Air Force Base. Air Force Office of
Special Investigations agent \textbf{Richard Doty} was assigned to him, and --- this is the
documented part, established through Greg Bishop's \emph{Project Beta} and by Doty's own later
statements --- Doty deliberately fed him fabricated material, including forged documents, in order
to discredit and misdirect him, and to steer him away from real classified activity at Kirtland.
Bennewitz's mental health deteriorated badly; he was eventually hospitalised, and he died in 2003.
The story continued without him.

The second source is \textbf{Thomas Castello}, allegedly a Dulce security officer, whose testimony
supplies the seven levels and Nightmare Hall. No such person has ever been documented to exist. No
employment record, no birth or death record matching the claimed details, no photographs of the man
--- and the ``Dulce Papers'' attributed to him are of unknown origin and appeared through third
parties.

The third source is \textbf{John Lear} and the 1987 \emph{Krill Papers} / Lear letter, which
assembled Dulce with MJ-12, the alien treaty and the abduction programme into the shape it holds
today, and which drew on Doty's fabrications. William Cooper then incorporated it into
\emph{Behold a Pale Horse} (1991), from which it entered general circulation.

So: the foundational documents are established forgeries produced by a US counterintelligence
officer to manipulate one distressed man, and the eyewitness cannot be shown to have existed.

And then the physical problems, which are decisive on their own. A seven-level facility of the
described scale requires excavation on the order of millions of cubic metres, which produces spoil
that must go somewhere visible; it requires power on the order of a small city, which requires
transmission lines or a generating plant; it requires ventilation shafts, cooling, water supply and
sewage; and it requires a workforce that must be housed, fed, paid and rotated. Dulce is a
community of about 2,\!600 people. There is no rail spur, no transmission corridor, no heavy road
traffic, and no visible spoil. Satellite imagery of Archuleta Mesa is publicly available at
sub-metre resolution and shows a mesa. The Jicarilla Apache Nation, which owns the land, has
repeatedly and flatly denied any such facility, and it is their land --- a point that gets
remarkably little attention from people who claim to be exposing the exploitation of the powerless.

Compare it to Cheyenne Mountain in Section~\ref{sec:cheyennemountain}, which is a real facility of comparable ambition,
and which required 700,\!000 tonnes of removed rock, a visible access road, published construction
contracts and five years of work by identifiable contractors. That is what building a hole in a
mountain looks like. It is not invisible. It cannot be.

\begin{funfact}
The Bennewitz affair is, to me, the darkest story in this book, and it has nothing to do with
aliens. A government agency identified a vulnerable civilian, exploited his beliefs with
deliberately manufactured evidence, and watched him disintegrate --- for operational convenience.
Doty has discussed it publicly. Nobody was disciplined. And the fabrications outlived the man by
decades and are still being repeated as testimony. If you want a reason to distrust official
sources, you do not need a saucer. You need a filing cabinet and a man with a badge and an
objective.
\end{funfact}

% ---------------------------------------------------------------
\section{Journey to The Center of The Earth}\label{sec:journeycenterearth}

The hollow Earth is worth a section because, like Section~\ref{sec:hollowmoon}, we can close it with measurements,
and because the measurements are genuinely interesting.

We know the interior structure of the Earth in detail, from seismology. Every large earthquake
illuminates the planet, and a global network of thousands of seismometers records the arrival times
of P-waves (compressional) and S-waves (shear) at every azimuth. From that you reconstruct the
interior the way a CT scanner reconstructs a body.

\figwrap[l]{0.42}{The Earth's interior as resolved by seismology: crust, mantle, liquid outer core, solid inner core.}{https://commons.wikimedia.org/wiki/File:Earth_poster.svg}{fig:earthinterior}
The results, established over a century: a crust averaging 35 km on continents and 7 km under
oceans, bounded by the Mohorovi\v{c}i\'c discontinuity found in 1909 (see Figure~\ref{fig:earthinterior}). A mantle to 2,\!890 km, with
transition-zone discontinuities at 410 and 660 km caused by known mineral phase changes ---
olivine to wadsleyite to ringwoodite to perovskite --- reproduced in diamond anvil cell
experiments at the predicted pressures. A liquid outer core from 2,\!890 to 5,\!150 km, identified
by Beno Gutenberg in 1913 from the S-wave shadow zone, because shear waves cannot propagate through
liquid. And a solid inner core, discovered by Inge Lehmann in 1936 from a faint reflected phase ---
one of the finest pieces of inference in the history of geophysics.

Independent confirmations: the planet's mean density of 5.51 g/cm$^3$ against surface rock at
around 2.7 requires a dense interior. The moment of inertia constrains the mass distribution. The
geodynamo requires a convecting conductive liquid layer. Free oscillation modes --- the whole planet
ringing after great earthquakes --- match the model. And the Kola Superdeep Borehole, drilled to
12.26 km over twenty years, was stopped because at 180$^\circ$C the rock behaved plastically and
the hole closed on the bit. That is 0.2\% of the way to the centre, and it was already too hot.

The pressure at the centre is about 360 gigapascals and the temperature is comparable to the
surface of the Sun. \hi{There is no cavity, and there could not be: rock has finite strength, and a
void of any scale collapses.}

\Needspace*{0.40\textheight}
\subsubsection*{The one witness the hollow-Earth literature cannot do without}

\figwrap[l]{0.40}{Byrd in flying kit before the Fokker F.VIIa/3m \emph{Josephine Ford}, the aircraft of the disputed 1926 North Pole flight.}{https://commons.wikimedia.org/wiki/File:Richard_E._Byrd_and_Fokker_F.VII_Josephine_Ford.jpg}{fig:josephineford}
There is a reason this section cannot end at the seismology, and his name is Richard Evelyn Byrd.
Every version of the modern hollow-Earth claim, without exception, is built on the assertion that a
United States Navy admiral flew into the polar opening and said so. Remove Byrd and the entire
edifice is Halley, Symmes and Verne again --- interesting history with no witness attached. So the
honest thing to do is to take the witness seriously, which means finding out who he actually was,
and that turns out to be far more interesting than the legend built on top of him.

The first thing to establish is that Byrd was not, in the phrase the field likes to use, ``just a
pilot'' (see Figure~\ref{fig:josephineford}). He was one of the most decorated and most famous
Americans alive. He held the Medal of Honor and the Navy Cross. He was promoted to rear admiral by
a special act of Congress on 21 December 1926, jumping the seniority list entirely, at the age of
thirty-eight. He was given a ticker-tape parade down Broadway. He was on the cover of \emph{Time}.
He led five Antarctic expeditions across three decades, ran the Navy's polar programme, and by the
end of his life had a mountain range, a research station and a coastal shelf named after him. When
he spoke about the poles, the United States press printed it verbatim, because he was the single
recognised authority on the subject in the English-speaking world.

That status is the whole point, and it is why he belongs in this book rather than in a footnote.
The pilot-credibility argument set out in Section~\ref{sec:pilots} is normally run on a captain
with twenty thousand hours. Byrd is that argument at maximum amplitude: the most credentialled
airman of his generation, a national hero, a flag officer, and a man whose profession was
precisely to observe and record the polar regions. \hi{If credentials alone settled questions of
fact, Byrd would settle this one.} He is the strongest form of the appeal to authority that this
subject has ever had available to it, which makes him the ideal place to test whether the appeal
works.

\Needspace*{0.40\textheight}
\subsubsection*{What Byrd actually flew, and what he actually claimed}
\figwrap[r]{0.42}{The Ford Trimotor \emph{Floyd Bennett}, lettered for the Byrd Antarctic Expedition --- the aircraft of the first flight over the South Pole, 28--29 November 1929.}{https://commons.wikimedia.org/wiki/File:Ford_Trimotor_Floyd_Bennett_of_the_Byrd_Antarctic_Expedition.jpg}{fig:floydbennett}
On 9 May 1926, flying the trimotor Fokker \emph{Josephine Ford} from Ny-{\AA}lesund in Spitsbergen
with Floyd Bennett at the controls, Byrd reported reaching the North Pole and returning in fifteen
hours and fifty-seven minutes. The claim made him famous, and it has been contested by aviation
historians ever since, on the arithmetic: the round trip is roughly 1,\!535 statute miles, and the
quoted time implies a ground speed the aircraft is not thought to have been capable of holding in
the prevailing conditions. In 1996 Byrd's original expedition diary was made available at Ohio
State University, and it contained erased-but-legible sextant sight reductions that do not match
the figures in his official report. Dennis Rawlins read that as proof he turned back short.
Gerald Newsom's subsequent analysis was more cautious and concluded the record cannot settle it
either way. What is not in dispute is that the primary navigational evidence is ambiguous, and
that it is ambiguous because of alterations made by Byrd himself.

The South Pole flight is a different matter (see Figure~\ref{fig:floydbennett}). On 28--29 November
1929, in the Ford Trimotor \emph{Floyd Bennett}, Byrd's crew flew from Little America over the pole
and back, dumping emergency rations to gain the altitude needed to clear the Liv Glacier. That
flight is corroborated, uncontested, and a genuine feat of airmanship. So the record of the most
famous polar aviator in history is: one flight that is certainly real and remarkable, and one that
is probably overstated. This is exactly the pattern Section~\ref{sec:pilots} predicts. Pilots are
reliable about procedure and instruments; the failure appears at the point where the claim depends
on a measurement nobody else was in a position to check.

It is worth being clear about what Byrd never claimed. In thirty years of published books,
lectures, official reports, radio broadcasts and congressional testimony, he never once described a
cavity, an opening, an inner sun, a subterranean civilisation, or a warm land at either pole. Not
once, in a body of primary material running to thousands of pages.

\Needspace*{0.40\textheight}
\subsubsection*{Operation Highjump: what the largest polar expedition in history was for}

\figwrap[r]{0.42}{Rear Admiral Byrd during Operation Highjump, working over the chart table --- the mapping mission that the legend later reinterpreted as a battle.}{https://commons.wikimedia.org/wiki/File:Rear_Admiral_Richard_Byrd,_Operation_Highjump.JPG}{fig:byrdhighjump}
The engine of the modern legend is Operation Highjump, and its real dimensions are impressive enough
that no invention is required (see Figure~\ref{fig:byrdhighjump}). Formally the United States Navy
Antarctic Developments Project, designated Task Force 68, it ran from August 1946 to February 1947
and deployed roughly 4,\!700 men, thirteen ships and thirty-three aircraft --- the largest expedition
ever sent to Antarctica, before or since. Rear Admiral Richard Cruzen was the operational commander;
Byrd was officer in charge of the expedition as a whole and flew with Task Group 68.4 from the
carrier USS \emph{Philippine Sea}. Four men died, three of them in the crash of a PBM Mariner on
Thurston Island.

What it was for is not a mystery, because the operation order says so and has been declassified for
decades: train personnel and test materiel in extreme cold, establish and maintain a base, extend
American sovereignty claims in practice, and survey the continent by air. It produced some
70,\!000 aerial photographs and mapped roughly 1.4~million square miles of coastline. The strategic
subtext was the Arctic, not Antarctica: the Navy expected its next war to be fought across the top
of the world against the Soviet Union, and Antarctica in the southern summer was the only place to
rehearse it at scale. That is also why the operation ended early --- the ice, the weather and the
lack of a second season's supply, not enemy action.

The legend requires Highjump to have been a combat operation defeated by a hidden power, and the
inventory makes that reading difficult to sustain (see Figure~\ref{fig:highjumpr4d}). Task Force 68
carried icebreakers, seaplane tenders, a submarine, tractors and photographic aircraft. It carried
no ground-attack force, no invasion troops, and no bombardment plan. Then it went home, published
its charts, and the Navy returned in 1947--48 for Operation Windmill to ground-truth the survey
control points, which is not what a service does after losing a battle. Section~\ref{sec:aryan}
deals with the Nazi-redoubt version of the same story and with why that particular myth exists.

\Needspace*{0.40\textheight}
\subsubsection*{The quotation that built an entire genre, and where it came from}
\figwrap[l]{0.42}{R4D Skytrains ranged on the flight deck of USS \emph{Philippine Sea} en route to Antarctica, 8 January 1947. Six were flown off the carrier to Little America to make the Highjump mapping flights.}{https://commons.wikimedia.org/wiki/File:USS_Philippine_Sea_(CV-47)_transporting_Douglas_R4D_aircraft_to_Antarctica,_8_January_1947_(80-G-K-7623).jpg}{fig:highjumpr4d}
Here is the crux, and it is a clean case study in how this field manufactures evidence. The famous
Byrd quotations that circulate in hollow-Earth books --- the lines about ``that enchanted continent
in the sky'' and about ``the land beyond the Pole'' being ``the center of the great unknown'' --- are
assembled from two sources of very different quality, and the assembly is the fraud.

Byrd did give an interview to the International News Service correspondent Lee van Atta aboard the
\emph{Mount Olympus}, published in the Santiago newspaper \emph{El Mercurio} on 5 March 1947. It is
real, and it is quoted more or less accurately as far as it goes.

\pullquote[Rear Admiral Richard E.\ Byrd, interviewed by Lee van Atta aboard USS \emph{Mount Olympus}, \emph{El Mercurio}, 5 March 1947]{The most important result of my observations and discoveries is the potential effect they have in relation to the security of the United States \ldots\ the necessity for the United States to maintain its position in Antarctica.}

What Byrd said, in context, is that the expedition's chief significance was strategic, that
Antarctica was the last unexplored region on Earth, and that a future war would be fought by
aircraft flying over the poles. It is a briefing about polar air power, delivered by the man who had
spent thirty years arguing for polar air power. There is no cavity in it, and no opening.

The hollow-Earth material comes from somewhere else entirely: F.\ Amadeo Giannini's \emph{Worlds
Beyond the Poles} (1959), which asserts a Byrd flight of 1,\!700 miles beyond the North Pole into
ice-free land with mountains, forests, rivers and living animals, and Raymond Bernard's \emph{The
Hollow Earth} (1964), which popularised Giannini and added the ``secret diary''. Neither cites a
primary document, because there is none. \hi{The flight described appears in no Navy record, no
logbook, no newspaper of the period and none of Byrd's own writings, and in the year it is supposed
to have occurred Byrd was not in the Arctic at all.} Colin Summerhayes and Peter Beeching traced the
entire construction in \emph{Polar Record} in 2007 and found the quotations to be either fabricated
outright or spliced from the \emph{El Mercurio} interview with the geography quietly changed.

The ``secret diary'' deserves its own sentence, because it is the load-bearing claim and it is the
easiest one to check.

\actualq[claim]{Byrd kept a hidden diary describing a warm land beyond the Pole, and it was
suppressed.}{Byrd's papers are open. They fill hundreds of linear feet at the Byrd Polar and
Climate Research Center at Ohio State University, they are catalogued and available to any
researcher, and they contain no such account --- while the one diary that genuinely was altered,
the 1926 North Pole flight diary, was altered by Byrd himself and argues against him rather than
for him.}

So the strongest witness the hollow Earth has is a man who never made the claim, in a document that
does not exist, about a flight that never happened. And the way to establish that was not to
dismiss him. It was to take him seriously enough to read him --- which is what he is owed. Byrd was
a superb organiser, a genuinely great expedition leader, a national hero on a scale no aviator
occupies today, and an unreliable narrator of his own instrument readings. All four of those are
true at once, and no amount of rank collapses them into one.

\begin{funfact}
The hollow Earth has a distinguished pedigree, which is why I treat it gently. Edmond Halley
proposed concentric shells in 1692 --- a serious attempt to explain geomagnetic secular variation,
using the best available physics. Leonhard Euler is often said to have supported it, probably
wrongly. John Cleves Symmes lobbied Congress in 1818 for an expedition to the polar opening, and
actually got a vote. Then Jules Verne in 1864 and Edgar Rice Burroughs's Pellucidar made it
literature, and the Nazi variant made it politics. Halley's version was good science that turned
out wrong. That is what good science looks like most of the time, and it is not an embarrassment.
\end{funfact}

\begin{srcnotes}
\item Summerhayes, C.\ P.\ \& Beeching, P.\ (2007). Hitler's Antarctic base: the myth and the reality.
\emph{Polar Record}, 43(1), 1--21 --- the definitive tracing of the Byrd quotations and of the
Antarctic-base literature to their actual sources.
\item United States Navy (1947). \emph{Report of Operation Highjump, U.S.\ Navy Antarctic
Developments Project 1947}, Task Force 68. Declassified; National Archives, Record Group 38.
\item Byrd, R.\ E.\ (1930). \emph{Little America: Aerial Exploration in the Antarctic and the Flight
to the South Pole}. New York: G.\ P.\ Putnam's Sons.
\item Byrd, R.\ E.\ (1938). \emph{Alone}. New York: G.\ P.\ Putnam's Sons.
\item Rawlins, D.\ (2000). Byrd's heroic 1926 North Pole failure. \emph{Polar Record}, 36(196),
25--50 --- the case against the 1926 claim, from the erased diary sights.
\item Newsom, G.\ H.\ (2002). Byrd's 1926 flight: the navigational evidence. Ohio State University
--- the more cautious reading, concluding the record is indeterminate.
\item Giannini, F.\ A.\ (1959). \emph{Worlds Beyond the Poles}. New York: Vantage Press --- the
origin of the ``land beyond the Pole'' attribution. Cited as the source of the claim, not as
evidence for it.
\item Bernard, R.\ (1964). \emph{The Hollow Earth}. New York: Fieldcrest --- popularised Giannini
and introduced the ``secret diary''.
\item Byrd Polar and Climate Research Center Archival Program, Ohio State University --- the Richard
E.\ Byrd Papers finding aid, showing the collection is catalogued and open.
\end{srcnotes}
